\documentclass[10pt,a4paper]{amsart}
\usepackage[margin=19mm]{geometry}
\usepackage{amsmath,amssymb,mathtools}
\usepackage[scr=rsfs,cal=euler]{mathalfa}
\usepackage{xfrac}
\usepackage[unicode,hidelinks,bookmarks=false]{hyperref}
\newcommand{\ie}{i.e.\ }
\newcommand{\cf}{cf.\ }
\newcommand{\resp}{respectively\ }
\newcommand{\Subsection}{\subsection}
\providecommand{\nobreakdash}{\nobreak\hbox{-}}
\setlength{\emergencystretch}{2em}
\multlinegap0pt
\usepackage{bm}
\usepackage{amssymb}
\usepackage{mathtools}
\usepackage{cancel}
\usepackage{xfrac}
\usepackage[shortlabels]{enumitem}
\usepackage{xcolor}
\usepackage{tikz}
\usepackage{float}
\newtheorem{lemma}{Lemma}
\newtheorem{corollary}{Corollary}
\title{Benign Nonconvexity of Synchronization Landscape Induced by Graph Skeletons}
\author{Hongjin Wu\textsuperscript{\normalfont *}, Ulrik Brandes\textsuperscript{\normalfont\textdagger}}
\date{Source: arXiv:2512.21935v3 (CC BY 4.0)}
\begin{document}
\maketitle

\noindent\emph{Source and licence.} Benign Nonconvexity of Synchronization Landscape Induced by Graph Skeletons. Version arXiv:2512.21935v3 (CC BY 4.0), distributed by arXiv under the Creative Commons Attribution 4.0 International licence, http://creativecommons.org/licenses/by/4.0/, as recorded in the arXiv licence metadata for this exact version and shown on the abstract page https://arxiv.org/abs/2512.21935v3. Full licence notice: \texttt{licenses/arxiv-2512.21935-CC-BY-4.0.txt}.\par\smallskip

\noindent\textit{Editorial scope.} Counted unit: the article's lemma lemma:leaflike (source0.tex lines 2180--2188). The article's complete original proof of it is delivered with it (source0.tex lines 2190--2330, one \texttt{proof} environment of 141 lines). No mathematics is written, repaired or re-derived: the statement and the proof are the article's own lines, in the article's own notation, and no retained line is substituted or re-flowed.  Retained uncounted as the source-local context the proof needs: lines 397--405; lines 697--699; lines 717--720; lines 1393--1396; lines 1838--1856; lines 2025--2040.  Nothing was omitted from any retained span, so the package carries no omission record.  No retained line contains a citation, so the wrapper carries no bibliography and no bibliography entry.  No retained line contains a figure environment, an image-inclusion command or drawing code, so no image is delivered and the package declares no asset dependency.  Editorial formatting only, in addition: an AMS \texttt{amsart} preamble that restates the article's own declarations and macros used by the retained lines, namely the environments \texttt{lemma}, \texttt{corollary} on separate counters, together with the standard packages those lines need.  The retained environments print in this document as Lemma 1, Corollary 1 in order of appearance, whereas the article prints the same items with the numbering of its own full document.  The complete native source of the article as distributed in the arXiv e-print for this exact version is bundled unchanged as \texttt{sources/191698/source0.tex}.
\begin{equation}\label{E_2}
\min_{\theta
\in\mathbb{R}^n}
E_G(\boldsymbol\theta)
=
\frac{1}{2}
\sum_{1\le i,j\le n}
A_{ij}\bigl(1-\cos(\theta_i-\theta_j)\bigr).
\end{equation}

\begin{equation}\label{formula:eq-ec}
\sum_{j\in N(i)} \boldsymbol v_j = \mu_i \boldsymbol v_i\quad\text{where}\quad \mu_i=\sum_{j\in N(i)}\cos(\theta_j-\theta_i),
\end{equation}

\begin{lemma}\label{stable-eq}
If $\boldsymbol{\theta}$ is a second-order stationary point of~\eqref{E_2},
then \eqref{formula:eq-ec} holds with $\mu_i\ge 0$ for every $1\le i\le n$.
\end{lemma}

\begin{corollary}\label{stable-closed-twins}
Let $\boldsymbol{\theta}$ be a second-order stationary point of $E_G$ in~\eqref{E_2}.
Then structural closed twins $a$ and $b$ must synchronize, i.e., $\theta_a=\theta_b \pmod{2\pi}$.
\end{corollary}

\begin{lemma}\label{pairwise-no-nauty-extra-node}
Let $\boldsymbol{\theta}$ be a second-order stationary point of $E_G$ on a graph
$G=(V,E)$.
Let $a,b\in V$ be adjacent nodes, and suppose that there exist disjoint sets
$S,T\subseteq V\setminus\{a,b\}$ such that
\[
N(a)\setminus\{b\}=S\sqcup T,
\qquad
N(b)\setminus\{a\}=S.
\]
If
\[
\theta_i=\theta_b \pmod{2\pi}, \qquad\text{for every } i\in T,
\]
then $a$ and $b$ synchronize, i.e.,
\[
\theta_a=\theta_b \pmod{2\pi}.
\]
\end{lemma}

\begin{lemma}\label{lem:opentwins}
Let $\boldsymbol\theta$ be a second-order stationary point of the energy
function~(2) of the Kuramoto model on a graph $G=(V,E)$.
Let $P,Q\subseteq V$ be disjoint nonempty sets,
and suppose that
\[
    P\subseteq N(i)\subseteq P\cup Q
    \qquad \text{for every } i\in Q.
\]
If the nodes in $Q$ are synchronized, with
$\mathbf{v}_i=\mathbf{v}$ for all $i\in Q$, then there exists
a scalar $\mu\geq 0$ such that
\[
    \sum_{j\in P}\mathbf{v}_j=\mu\mathbf{v}.
\]
\end{lemma}

\begin{lemma}[Propagation of the leaf-like property]
\label{lemma:leaflike}
Let $\boldsymbol{\theta}\in\mathbb{R}^{|V|}$ be a second-order
stationary point of the Kuramoto energy $E_G(\boldsymbol{\theta})$ on a connected
quasi-threshold graph $G=(V,E)$ with underlying rooted tree $T$.
If $a\in V(T)$ is a non-leaf node and all its children are
leaf-like in $T$ at $\boldsymbol{\theta}$, then $a$ is also
leaf-like in $T$ at $\boldsymbol{\theta}$.
\end{lemma}

\begin{proof}[Proof of Lemma~\ref{lemma:leaflike}]
If $a$ has exactly one child $b$, then
\[
    N_G[a]=N_G[b]
    =\mathrm{Anc}_T(a)\cup\{a,b\}\cup\mathrm{Desc}_T(b).
\]
Thus, $a$ and $b$ are structural closed twins in $G$ and
synchronize at $\boldsymbol{\theta}$ by
Corollary~\ref{stable-closed-twins}.
Since $b$ is leaf-like in $T$ at $\boldsymbol{\theta}$ and
\[
    \mathrm{Desc}_T(a)=\{b\}\cup\mathrm{Desc}_T(b),
\]
every descendant of $a$ synchronizes with $a$.
Hence, $a$ is leaf-like in $T$ at $\boldsymbol{\theta}$.
It remains to consider the case where $a$ has at least two
children.

\medskip
\noindent
\textit{Step 1: Applying the synchronous extension lemma.}
Set
\[
    P:=\{a\}\cup\mathrm{Anc}_T(a),
    \qquad
    \mathbf{s}:=\sum_{j\in P}\mathbf{v}_j
    =\mathbf{v}_a+\sum_{j\in\mathrm{Anc}_T(a)}\mathbf{v}_j.
\]
Fix a child $i$ of $a$ and let
\[
    Q_i:=\{i\}\cup\mathrm{Desc}_T(i).
\]
Since $i$ is leaf-like in $T$ at $\boldsymbol{\theta}$,
we have $\mathbf{v}_k=\mathbf{v}_i$ for every $k\in Q_i$.
Moreover, adjacency in $G$ is determined by the
ancestor--descendant relation in $T$.
Thus, every node $k\in Q_i$ is adjacent to every node in $P$
and has no neighbors outside $P\cup Q_i$; that is,
\[
    P\subseteq N_G(k)\subseteq P\cup Q_i
    \qquad \text{for every } k\in Q_i.
\]
The sets $P$ and $Q_i$ are disjoint and nonempty.
Applying Lemma~\ref{lem:opentwins} therefore gives a scalar
$\alpha_i\ge0$ such that
\[
    \mathbf{s}=\alpha_i\mathbf{v}_i.
\]
Consequently, if $\mathbf{s}\neq\boldsymbol{0}$, then
$\alpha_i=\|\mathbf{s}\|>0$ and
\[
    \mathbf{v}_i=\frac{\mathbf{s}}{\|\mathbf{s}\|}
    \qquad \text{for every } i\in\mathrm{Children}_T(a).
\]
Hence all children of $a$ synchronize whenever
$\mathbf{s}\neq\boldsymbol{0}$.

\medskip
\noindent
\textit{Step 2: Excluding a vanishing external sum.}
Suppose, for contradiction, that $\mathbf{s}=\boldsymbol{0}$.
Write
\[
    \mathbf{u}:=\sum_{j\in\mathrm{Anc}_T(a)}\mathbf{v}_j,
    \qquad
    \mathbf{d}:=\sum_{j\in\mathrm{Desc}_T(a)}\mathbf{v}_j.
\]
Then $\mathbf{u}=-\mathbf{v}_a$.
Since
\[
    N_G(a)=\mathrm{Anc}_T(a)\cup\mathrm{Desc}_T(a),
\]
Lemma~\ref{stable-eq} gives
\[
    \mathbf{u}+\mathbf{d}=\lambda\mathbf{v}_a
    \qquad \text{for some } \lambda\ge0.
\]
It follows that
\[
    \mathbf{d}=(\lambda+1)\mathbf{v}_a.
\]

Let $H=\nabla^2 E_G(\boldsymbol{\theta})$ and define
$\boldsymbol{x}\in\mathbb{R}^{|V|}$ by
\[
    x_k=
    \begin{cases}
        1, & k\in\{a\}\cup\mathrm{Desc}_T(a),\\
        0, & \text{otherwise}.
    \end{cases}
\]
By the ancestor--descendant structure of $G$, the edges
crossing from $\{a\}\cup\mathrm{Desc}_T(a)$ to its complement
are precisely all edges joining this set to
$\mathrm{Anc}_T(a)$.
Therefore,
\[
\begin{aligned}
    \boldsymbol{x}^{\top}H\boldsymbol{x}
    &=
    \sum_{k\in\{a\}\cup\mathrm{Desc}_T(a)}
    \sum_{j\in\mathrm{Anc}_T(a)}
    \cos(\theta_k-\theta_j)\\
    &=\langle\mathbf{v}_a+\mathbf{d},\mathbf{u}\rangle\\
    &=\langle(\lambda+2)\mathbf{v}_a,-\mathbf{v}_a\rangle\\
    &=-(\lambda+2)<0,
\end{aligned}
\]
where we used $\|\mathbf{v}_a\|=1$.
This contradicts the positive semidefiniteness of $H$.
Thus $\mathbf{s}\neq\boldsymbol{0}$, and Step~1 implies
that all children of $a$ synchronize.
Since each child is leaf-like in $T$ at
$\boldsymbol{\theta}$, all descendants of $a$ synchronize
with one another.

\medskip
\noindent
\textit{Step 3: Synchronizing $a$ with its descendants.}
Fix a child $b$ of $a$ and set
\[
\begin{aligned}
    S&:=\mathrm{Anc}_T(a)\cup\mathrm{Desc}_T(b),\\
    R&:=\mathrm{Desc}_T(a)\setminus
    \bigl(\{b\}\cup\mathrm{Desc}_T(b)\bigr).
\end{aligned}
\]
These are disjoint subsets of $V\setminus\{a,b\}$ satisfying
\[
    N_G(a)\setminus\{b\}=S\sqcup R,
    \qquad
    N_G(b)\setminus\{a\}=S.
\]
Every node in $R$ synchronizes with $b$, because all
descendants of $a$ are synchronized.
Applying Lemma~\ref{pairwise-no-nauty-extra-node} to the
adjacent nodes $a,b$, with common neighbors $S$ and extra
neighbors $R$, yields $\mathbf{v}_a=\mathbf{v}_b$.
Thus every descendant of $a$ synchronizes with $a$,
so $a$ is leaf-like in $T$ at $\boldsymbol{\theta}$.
\end{proof}

\end{document}
